The \code{BPatch\_edge} class represents a control flow edge in a
\code{BPatch\_flowGraph}.

\begin{apient}
  BPatch_point *getPoint()
\end{apient}

\apidesc{
  Return an instrumentation point for this edge. This point can be
  passed to \code{BPatch\_process::insertSnippet} to
  instrument the edge.
}

\begin{apient}
  enum BPatch_edgeType {
    CondJumpTaken,
    CondJumpNottaken,
    UncondJump,
    NonJump
  }
\end{apient}

\apidesc{}

\begin{apient}
  BPatch_edgeType getType()
\end{apient}

\apidesc{
  Return a type describing this edge. A CondJumpTaken edge is found
  after a conditional branch, along the edge that is
  taken when the condition is true. A CondJumpNottaken edge follows
  the path when the condition is not taken.
  UncondJump is used along an edge that flows out of an unconditional
  branch that is always taken. NonJump is an edge
  that flows out of a basic block that does not end in a jump, but
  falls through into the next basic block.
}
\begin{apient}
  BPatch_basicBlock *getSource()
\end{apient}

\apidesc{
  Return the source \code{BPatch\_basicBlock} that this edge flows from.
}
\begin{apient}
  BPatch_basicBlock *getTarget()
\end{apient}

\apidesc{
  Return the target \code{BPatch\_basicBlock} that this edge flows to.
}
\begin{apient}
  BPatch_flowGraph *getFlowGraph()
\end{apient}

\apidesc{
  Returns the CFG that contains the edge.
}
